\documentclass[en,hazy,blue,12pt,normal]{elegantnote}

\input{preface}

\title{Winter 2026 DSC 240: Homework 0}
\author{Zeyu Bian}
\date{\today}

\begin{document}

\maketitle

\begin{homework}
 Refreshers on Optimization and probability fundamentals.
\end{homework}

\begin{enumerate}[label=\textbf{(\alph*)}, leftmargin=2.2em]

\item \textbf{(Continuous optimization)}
Let $x_1,\dots,x_n \in \mathbb{R}$ and
\[
f(\theta)=\sum_{i=1}^n w_i(x_i-\theta)^2.
\]
Assume first that $w_i>0$ for all $i$. Expand and differentiate:
\[
f(\theta)=\sum_{i=1}^n w_i(x_i^2-2x_i\theta+\theta^2)
= \Big(\sum_{i=1}^n w_i\Big)\theta^2 - 2\Big(\sum_{i=1}^n w_i x_i\Big)\theta + \sum_{i=1}^n w_i x_i^2.
\]
Then
\[
f'(\theta)=2\Big(\sum_{i=1}^n w_i\Big)\theta - 2\sum_{i=1}^n w_i x_i.
\]
Setting $f'(\theta)=0$ gives
\[
\theta^\ast=\frac{\sum_{i=1}^n w_i x_i}{\sum_{i=1}^n w_i}.
\]
Moreover, when all $w_i>0$, $\sum_i w_i>0$ so $f$ is a strictly convex quadratic and $\theta^\ast$ is the unique global minimizer.

\medskip
\noindent \textbf{If some $w_i$ are negative:}
$f(\theta)$ is still a quadratic in $\theta$ with leading coefficient $a=\sum_{i=1}^n w_i$.
\begin{itemize}
  \item If $\sum_i w_i>0$, then $f$ still opens upward and the (unique) minimizer is still
  $\displaystyle \theta^\ast=\frac{\sum_i w_i x_i}{\sum_i w_i}$.
  \item If $\sum_i w_i=0$, then $f(\theta)$ becomes linear in $\theta$ (unless also $\sum_i w_i x_i=0$, in which case it is constant). In the generic linear case it has no minimizer over $\mathbb{R}$.
  \item If $\sum_i w_i<0$, then $f$ opens downward and is unbounded below, so there is no minimizer.
\end{itemize}

\item \textbf{(Counting and combinatorics)}
There are $2n$ kids, split uniformly at random into two groups of size $n$. Consider selecting Group 1 uniformly among all $\binom{2n}{n}$ size-$n$ subsets.

Let the two tallest kids be $A$ and $B$.
\begin{itemize}
\item \emph{Same subgroup:} Either both in Group 1 or both in Group 2.
  \[
  \mathbb{P}(\text{both in Group 1})=\frac{\binom{2n-2}{n-2}}{\binom{2n}{n}}, \qquad
  \mathbb{P}(\text{both in Group 2})=\frac{\binom{2n-2}{n}}{\binom{2n}{n}}.
  \]
  Since $\binom{2n-2}{n}=\binom{2n-2}{n-2}$, we get
  \[
  \mathbb{P}(\text{same})=
  \frac{2\binom{2n-2}{n-2}}{\binom{2n}{n}}
  =\frac{n-1}{2n-1}.
  \]
\item \emph{Different subgroups:}
  \[
  \mathbb{P}(\text{different})=1-\mathbb{P}(\text{same})=1-\frac{n-1}{2n-1}=\frac{n}{2n-1}.
  \]
\end{itemize}
\emph{Check for $n=2$:} $\mathbb{P}(\text{same})=\frac{1}{3}$ and $\mathbb{P}(\text{different})=\frac{2}{3}$, which matches a manual enumeration for 4 kids.

\item \textbf{(Bayes rule)}
Let $K$ be the event ``candidate knows the answer'' and $C$ be the event ``candidate answers correctly.'' Given:
\[
\mathbb{P}(K)=p,\quad \mathbb{P}(C\mid K)=0.99,\quad \mathbb{P}(C\mid K^c)=\frac{1}{k}.
\]
First compute $\mathbb{P}(C)$:
\[
\mathbb{P}(C)=\mathbb{P}(C\mid K)\mathbb{P}(K)+\mathbb{P}(C\mid K^c)\mathbb{P}(K^c)
=0.99p+\frac{1}{k}(1-p).
\]
By Bayes' rule,
\[
\mathbb{P}(K\mid C)=\frac{\mathbb{P}(C\mid K)\mathbb{P}(K)}{\mathbb{P}(C)}
=\frac{0.99\,p}{0.99\,p+\frac{1}{k}(1-p)}.
\]

\item \textbf{(Likelihood and maximum likelihood)}
We have likelihood
\[
L(p)=p^6(1-p)^4,\qquad 0<p<1.
\]
Maximizing $\log L(p)$ preserves the $\arg\max$ because $\log(\cdot)$ is strictly increasing:
\[
\arg\max_p L(p) \;=\;\arg\max_p \log L(p).
\]
Now
\[
\ell(p):=\log L(p)=6\log p+4\log(1-p).
\]
Differentiate and set to zero:
\[
\ell'(p)=\frac{6}{p}-\frac{4}{1-p}=0
\;\Longrightarrow\;
6(1-p)=4p
\;\Longrightarrow\;
6=10p
\;\Longrightarrow\;
p^\ast=\frac{6}{10}=0.6.
\]
Also $\ell''(p)=-\frac{6}{p^2}-\frac{4}{(1-p)^2}<0$, so this is indeed the (unique) maximizer.

\item \textbf{(Calculus, gradients)}
Let $w\in\mathbb{R}^d$, $x_i\in\mathbb{R}^d$ (column vectors), $y_i\in\mathbb{R}$, and $\lambda\ge 0$.
\[
F(w)=\sum_{i=1}^n (x_i^T w-y_i)^2 + \lambda\sum_{j=1}^d w_j^2.
\]
For each $i$, define $r_i(w)=x_i^T w-y_i$ (a scalar). Then $\nabla_w r_i(w)=x_i$. By the chain rule,
\[
\nabla_w \big(r_i(w)^2\big)=2r_i(w)\nabla_w r_i(w)=2(x_i^T w-y_i)x_i.
\]
Also $\nabla_w\big(\lambda \sum_{j=1}^d w_j^2\big)=2\lambda w$. Therefore,
\[
\nabla F(w)=2\sum_{i=1}^n (x_i^T w-y_i)x_i + 2\lambda w.
\]
\textbf{Matrix form (optional):} If $X\in\mathbb{R}^{n\times d}$ has rows $x_i^T$ and $y\in\mathbb{R}^n$ has entries $y_i$, then
\[
F(w)=\|Xw-y\|_2^2+\lambda\|w\|_2^2,
\qquad
\nabla F(w)=2X^T(Xw-y)+2\lambda w.
\]

\item \textbf{(Chain-rule and softmax / log-sum-exp)}
Let
\[
f(x_1,\dots,x_n)=\log\left(\sum_{i=1}^n e^{x_i}\right).
\]
Let $S=\sum_{i=1}^n e^{x_i}$. Then
\[
\frac{\partial f}{\partial x_j}
=\frac{1}{S}\cdot \frac{\partial}{\partial x_j}\left(\sum_{i=1}^n e^{x_i}\right)
=\frac{1}{S}\, e^{x_j}.
\]
Thus the gradient is
\[
\nabla f(x)=\begin{bmatrix}
\frac{e^{x_1}}{\sum_{i=1}^n e^{x_i}}\\[4pt]
\vdots\\[2pt]
\frac{e^{x_n}}{\sum_{i=1}^n e^{x_i}}
\end{bmatrix},
\]
which is exactly the \emph{softmax} probability vector.

\item \textbf{(Simple mathematical proof)}
Let $m=\max_i x_i$.

\medskip
\noindent \textbf{Lower bound:} Since $e^{x_i}\ge 0$ and one term attains $e^{m}$,
\[
\sum_{i=1}^n e^{x_i}\ge e^{m}
\quad\Longrightarrow\quad
f(x)=\log\left(\sum_{i=1}^n e^{x_i}\right)\ge \log(e^{m})=m.
\]

\medskip
\noindent \textbf{Upper bound:} Since $x_i\le m$ for all $i$, we have $e^{x_i}\le e^{m}$ for all $i$. Therefore,
\[
\sum_{i=1}^n e^{x_i}\le \sum_{i=1}^n e^{m}=n e^{m}
\quad\Longrightarrow\quad
f(x)=\log\left(\sum_{i=1}^n e^{x_i}\right)\le \log(n e^{m})=m+\log n.
\]

\medskip
\noindent Combining both inequalities yields
\[
\max_i x_i \;\le\; f(x_1,\dots,x_n)\;\le\;\max_i x_i+\log n.
\]

\end{enumerate}

\begin{homework}
     Refreshers on time complexity, data structure and python / numpy. Please write codes for python3.
\end{homework}

\begin{enumerate}[label=\textbf{(\alph*)}, leftmargin=2.2em]

\item \textbf{(Numpy matrix creation and multiplication)}
Construct:
\[
A\in\mathbb{R}^{5\times 4}=\begin{bmatrix}
1&2&3&4\\
5&6&7&8\\
9&10&11&12\\
13&14&15&16\\
17&18&19&20
\end{bmatrix},
\qquad
B\in\mathbb{R}^{4\times 3}=\begin{bmatrix}
1&5&9\\
2&6&10\\
3&7&11\\
4&8&12
\end{bmatrix},
\]
where $B$ is filled \emph{top-to-bottom, then left-to-right} (column-major order).

\begin{lstlisting}[style=py]
import numpy as np

def make_A_B():
    # A: 1..20 left-to-right then top-to-bottom (row-major fill)
    A = np.arange(1, 21).reshape(5, 4)

    # B: 1..12 top-to-bottom then left-to-right (column-major fill)
    B = np.arange(1, 13).reshape(4, 3, order="F")
    return A, B

def matmul(A, B):
    """Return the matrix product AB."""
    return A @ B   # equivalently: np.dot(A, B)

# Example:
# A, B = make_A_B()
# C = matmul(A, B)
\end{lstlisting}

\item \textbf{(Sparse matrix-vector multiplication)}
\textbf{Dense case:} Multiplying an $n\times n$ dense matrix by a dense vector costs
\[
O(n^2)
\]
time in the worst case (because there are $n^2$ multiply-add contributions).

\medskip
\noindent \textbf{Sparse case:} If $A$ is sparse with $\mathrm{nnz}(A)$ nonzero entries, sparse matrix-vector multiplication costs
\[
O(\mathrm{nnz}(A))
\]
time (up to constant factors depending on storage format), since you only process nonzero entries.

\medskip
\noindent \textbf{Construct $A\in\mathbb{R}^{(n-1)\times n}$ so that $Ax=[x_1-x_2,\dots,x_{n-1}-x_n]^T$:}
This is the first-difference operator:
\[
A=\begin{bmatrix}
1 & -1 & 0 & \cdots & 0\\
0 & 1 & -1 & \cdots & 0\\
\vdots & & \ddots & \ddots & \vdots\\
0 & \cdots & 0 & 1 & -1
\end{bmatrix}.
\]

\begin{lstlisting}[style=py]
import numpy as np
from scipy import sparse

def difference_matrix(n: int):
    """
    Return a sparse matrix A of shape ((n-1), n) such that
    for any x in R^n, A @ x = [x1-x2, x2-x3, ..., x_{n-1}-x_n]^T.
    """
    if n < 2:
        raise ValueError("n must be >= 2 so that (n-1) x n is non-empty.")

    data = np.ones(n - 1)
    # diagonal 0 has +1's, diagonal +1 has -1's
    A = sparse.diags([data, -data], offsets=[0, 1], shape=(n - 1, n), format="csr")
    return A

# Example check:
# n = 5
# A = difference_matrix(n)
# x = np.array([1, 3, 6, 10, 15])
# print(A @ x)  # should be [1-3, 3-6, 6-10, 10-15] = [-2, -3, -4, -5]
\end{lstlisting}

\item \textbf{(Manipulating text using python, data structures)}
We can use a \textbf{dictionary} (\texttt{dict}) mapping each word to its count. For basic tokenization, one common approach is to lowercase and extract alphanumeric ``words'' using a regex.

\begin{lstlisting}[style=py]
import re

def word_counts(filepath: str):
    """
    Read a text file line by line and return a dictionary:
    {word -> count}.
    Words are case-folded to lowercase and extracted via regex.
    """
    counts = {}  # dict: word -> int count

    with open(filepath, "r", encoding="utf-8") as f:
        for line in f:
            # Extract tokens (letters/digits/underscore). You may adjust pattern if desired.
            words = re.findall(r"\w+", line.lower())
            for w in words:
                counts[w] = counts.get(w, 0) + 1

    return counts

# Example:
# counts = word_counts("data_example.txt")
# print(len(counts), "unique words")
# print(counts.get("machine", 0))
\end{lstlisting}

\noindent \textbf{Data structures used:} A \texttt{dict} for counts (hash table), plus temporary \texttt{list} of tokens per line from \texttt{re.findall}.

\medskip
\noindent \textbf{Time complexity:} Let $T$ be the total number of word tokens extracted from the file. Each dictionary increment is $O(1)$ expected time, so total time is $O(T)$ (plus the linear pass to read and tokenize the text, which is also linear in the input size).

\medskip
\noindent \textbf{Space complexity:} Let $U$ be the number of unique words. The dictionary stores $U$ keys, so space is $O(U)$ (ignoring the input file storage, and treating tokenization overhead per line as transient).

\item \textbf{(Recursion with bookkeeping in python)}
Define Fibonacci numbers with $F_0=0$, $F_1=1$, and $F_n=F_{n-1}+F_{n-2}$ for $n\ge 2$.
A memoized recursion computes each $F_m$ at most once.

\begin{lstlisting}[style=py]
# Global memo dictionary for bookkeeping:
_fib_memo = {0: 0, 1: 1}

def fibonacci(n: int) -> int:
    """
    Return the nth Fibonacci number using recursion + memoization.
    Uses convention: F_0=0, F_1=1.
    """
    if n < 0:
        raise ValueError("n must be nonnegative.")
    if n in _fib_memo:
        return _fib_memo[n]
    _fib_memo[n] = fibonacci(n - 1) + fibonacci(n - 2)
    return _fib_memo[n]

# Example:
# print(fibonacci(10))  # 55
\end{lstlisting}

\noindent \textbf{Time complexity (with bookkeeping):} $O(n)$, since each value $F_0,\dots,F_n$ is computed once and stored.

\medskip
\noindent \textbf{Space complexity:} $O(n)$ for the memo table (and also $O(n)$ worst-case recursion depth).

\end{enumerate}

\end{document}